\section{Firmware}
\label{sec:firmware}

O firmware foi escrito em C sobre o ESP-IDF v6 \citep{espidf}, o \textit{framework} oficial da Espressif, que inclui o FreeRTOS \citep{freertos}, a pilha TCP/IP, os drivers de periféricos e bibliotecas de cliente/servidor HTTP. Um esboço inicial em Arduino (\texttt{esp32-petcare.ino}) foi mantido no repositório como referência, mas o firmware usado no protótipo é o ESP-IDF, organizado nos módulos da Tabela~\ref{tab:modulos}.

\begin{table}[!ht]
\caption{Módulos do firmware ESP-IDF (\texttt{firmware/esp32-petcare/main/}).}
\label{tab:modulos}
\centering
\small
\begin{tblr}{
  colspec={X[1,l] X[2,l]},
  row{1}={font=\bfseries},
  hline{1,Z}={0.8pt}, hline{2}={0.4pt},
  rowsep=1pt
}
Módulo & Responsabilidade \\
\texttt{main.c} & inicialização e laço principal (botão, Wi-Fi e sonar) \\
\texttt{config\_store} & configuração persistente na NVS; botão BOOT por 3\,s apaga e reinicia \\
\texttt{wifi\_manager} & modo estação, ponto de acesso \texttt{PetCare-Setup} e servidor HTTP de pareamento \\
\texttt{api\_client} & fila de leituras, \textit{heartbeat}, descoberta UDP e cliente HTTP \\
\texttt{sonar\_sensor} & UART2, analisador de pacotes e detecção de entrada/saída \\
\texttt{mic\_sensor} & ADC contínuo, análise por quadros e detecção de sons \\
\texttt{json\_util} & escape de texto interpolado em JSON \\
\texttt{capacitive\_sensor} & leitura do CR18-8DP (somente no protótipo; \pendente{publicar}) \\
\end{tblr}
\end{table}

\subsection{Concorrência}

O laço principal de \texttt{app\_main} verifica periodicamente o botão BOOT, a máquina de estados do Wi-Fi e a UART do sonar. O microfone roda em uma \textit{task} própria fixada no núcleo~1, já que a pilha Wi-Fi do ESP32 opera preferencialmente no núcleo~0, de modo que a análise de áudio não concorre com a rede. O envio para a API é feito por uma terceira \textit{task} (\texttt{api\_worker}), que consome uma fila FreeRTOS com até 16 leituras; os sensores apenas enfileiram, sem bloquear à espera da rede.

\subsection{Configuração persistente e pareamento}

SSID, senha, URL da API, \textit{token} do dispositivo e prefixo da placa (padrão \texttt{esp32-01}) ficam na partição NVS \citep{espidfnvs}, no \textit{namespace} \texttt{petcare\_cfg}. Nenhuma credencial é gravada no código.

Sem rede salva, ou após 10\,s sem conseguir se conectar, a placa abre o ponto de acesso aberto \texttt{PetCare-Setup} em \texttt{10.73.73.1}, comportamento confirmado no registro serial do protótipo. Nesse modo, um servidor HTTP embarcado expõe \texttt{/status}, \texttt{/scan} (redes visíveis) e \texttt{/provision} (recebe SSID, senha e \textit{token}), com cabeçalhos CORS para que a página de configurações da aplicação web converse diretamente com a placa. Se já houver rede salva, a placa tenta reconectar a cada 30\,s enquanto ninguém estiver conectado ao ponto de acesso, e o fecha quando a rede volta (por exemplo, após o roteador reiniciar depois de uma queda de energia). A sub-rede 10.73.73.0/24 foi escolhida para não colidir com redes domésticas comuns.

\subsection{Descoberta do servidor e envio}

Com a URL da API vazia, ou após uma falha de conexão, a placa envia em \textit{broadcast} UDP, na porta 3333, a mensagem \texttt{PETCARE\_DISCOVER}. A API responde \texttt{PETCARE\_API:<porta>}, e a placa monta a URL a partir do IP de origem da resposta e a salva na NVS. Assim, a placa encontra o servidor mesmo que o IP do computador mude por DHCP.

O envio usa o cliente HTTP do ESP-IDF com \textit{timeout} de 4\,s e o cabeçalho \texttt{x-device-token} quando configurado. Um \textit{heartbeat} (\texttt{POST /pairing/heartbeat}, com prefixo, SSID, IP e MAC) é enviado logo após a conexão e a cada 15\,s, com novas tentativas a cada 3\,s em caso de falha. Os textos interpolados no JSON (SSID, prefixo) são escapados.

\subsection{Limitações do firmware}

\begin{itemize}
  \item Leituras que falham no envio não são reenfileiradas, e o enfileiramento não espera quando a fila está cheia: \textbf{eventos ocorridos com a rede fora do ar são perdidos}. Um \textit{buffer} persistente com reenvio seria necessário para não perder visitas à caixa de areia.
  \item As leituras não levam horário da placa: o servidor registra o horário de chegada (o contrato aceita \texttt{recordedAt}, mas o firmware ainda não o preenche, pois não há sincronização de relógio).
  \item Os limiares de detecção (sonar e microfone) são constantes de compilação, ainda sem calibração pela interface.
\end{itemize}
